\documentclass[11pt,addpoints,answers]{exam}
\usepackage[margin=1in]{geometry}
\usepackage{amsmath, amsfonts}
\usepackage{enumerate}
\usepackage{graphicx}
\usepackage{titling}
\usepackage{url}
\usepackage{xfrac}
\usepackage{geometry}
\usepackage{graphicx}
\usepackage{natbib}
\usepackage{amsmath}
\usepackage{amssymb}
\usepackage{amsthm}
\usepackage{paralist}
\usepackage{epstopdf}
\usepackage{tabularx}
\usepackage{longtable}
\usepackage{multirow}
\usepackage{multicol}
\usepackage[colorlinks=true,urlcolor=blue]{hyperref}
\usepackage{fancyvrb}
\usepackage{algorithm}
\usepackage{algorithmic}
\usepackage{float}
\usepackage{paralist}
\usepackage[svgname]{xcolor}
\usepackage{enumerate}
\usepackage{array}
\usepackage{times}
\usepackage{url}
\usepackage{comment}
\usepackage{environ}
\usepackage{times}
\usepackage{textcomp}
\usepackage{caption}
\usepackage[colorlinks=true,urlcolor=blue]{hyperref}
\usepackage{listings}
\usepackage{parskip} % For NIPS style paragraphs.
\usepackage[compact]{titlesec} % Less whitespace around titles
\usepackage[inline]{enumitem} % For inline enumerate* and itemize*
\usepackage{datetime}
\usepackage{comment}
% \usepackage{minted}
\usepackage{lastpage}
\usepackage{color}
\usepackage{xcolor}
\usepackage{listings}
\usepackage{tikz}
\usetikzlibrary{shapes,decorations,bayesnet}
%\usepackage{framed}
\usepackage{booktabs}
\usepackage{cprotect}
\usepackage{xcolor}
\usepackage{verbatimbox}
\usepackage[many]{tcolorbox}
\usepackage{cancel}
\usepackage{wasysym}
\usepackage{mdframed}
\usepackage{subcaption}
\usetikzlibrary{shapes.geometric}

%%%%%%%%%%%%%%%%%%%%%%%%%%%%%%%%%%%%%%%%%%%
% Formatting for \CorrectChoice of "exam" %
%%%%%%%%%%%%%%%%%%%%%%%%%%%%%%%%%%%%%%%%%%%

\CorrectChoiceEmphasis{}
\checkedchar{\blackcircle}

%%%%%%%%%%%%%%%%%%%%%%%%%%%%%%%%%%%%%%%%%%%
% Better numbering                        %
%%%%%%%%%%%%%%%%%%%%%%%%%%%%%%%%%%%%%%%%%%%

\numberwithin{equation}{section} % Number equations within sections (i.e. 1.1, 1.2, 2.1, 2.2 instead of 1, 2, 3, 4)
\numberwithin{figure}{section} % Number figures within sections (i.e. 1.1, 1.2, 2.1, 2.2 instead of 1, 2, 3, 4)
\numberwithin{table}{section} % Number tables within sections (i.e. 1.1, 1.2, 2.1, 2.2 instead of 1, 2, 3, 4)


%%%%%%%%%%%%%%%%%%%%%%%%%%%%%%%%%%%%%%%%%%%
% Common Math Commands                    %
%%%%%%%%%%%%%%%%%%%%%%%%%%%%%%%%%%%%%%%%%%%
\input{mathabbreviations.tex}

%%%%%%%%%%%%%%%%%%%%%%%%%%%%%%%%%%%%%%%%%%%
% Code highlighting with listings         %
%%%%%%%%%%%%%%%%%%%%%%%%%%%%%%%%%%%%%%%%%%%

\definecolor{bluekeywords}{rgb}{0.13,0.13,1}
\definecolor{greencomments}{rgb}{0,0.5,0}
\definecolor{redstrings}{rgb}{0.9,0,0}
\definecolor{light-gray}{gray}{0.95}

\newcommand{\MYhref}[3][blue]{\href{#2}{\color{#1}{#3}}}%

\definecolor{dkgreen}{rgb}{0,0.6,0}
\definecolor{gray}{rgb}{0.5,0.5,0.5}
\definecolor{mauve}{rgb}{0.58,0,0.82}

\lstdefinelanguage{Shell}{
  keywords={tar, cd, make},
  %keywordstyle=\color{bluekeywords}\bfseries,
  alsoletter={+},
  ndkeywords={python, py, javac, java, gcc, c, g++, cpp, .txt, octave, m, .tar},
  %ndkeywordstyle=\color{bluekeywords}\bfseries,
  identifierstyle=\color{black},
  sensitive=false,
  comment=[l]{//},
  morecomment=[s]{/*}{*/},
  commentstyle=\color{purple}\ttfamily,
  stringstyle=\color{red}\ttfamily,
  morestring=[b]',
  morestring=[b]",
  backgroundcolor = \color{light-gray}
}

\lstset{columns=fixed, basicstyle=\ttfamily,
    backgroundcolor=\color{light-gray},xleftmargin=0.5cm,frame=tlbr,framesep=4pt,framerule=0pt}



%%%%%%%%%%%%%%%%%%%%%%%%%%%%%%%%%%%%%%%%%%%
% Custom box for highlights               %
%%%%%%%%%%%%%%%%%%%%%%%%%%%%%%%%%%%%%%%%%%%

% Define box and box title style
\tikzstyle{mybox} = [fill=blue!10, very thick,
    rectangle, rounded corners, inner sep=1em, inner ysep=1em]

% \newcommand{\notebox}[1]{
% \begin{tikzpicture}
% \node [mybox] (box){%
%     \begin{minipage}{\textwidth}
%     #1
%     \end{minipage}
% };
% \end{tikzpicture}%
% }

\NewEnviron{notebox}{
\begin{tikzpicture}
\node [mybox] (box){
    \begin{minipage}{\textwidth}
        \BODY
    \end{minipage}
};
\end{tikzpicture}
}

%%%%%%%%%%%%%%%%%%%%%%%%%%%%%%%%%%%%%%%%%%%
% Commands showing / hiding solutions     %
%%%%%%%%%%%%%%%%%%%%%%%%%%%%%%%%%%%%%%%%%%%

%% To HIDE SOLUTIONS (to post at the website for students), set this value to 0: \def\issoln{0}
\def\issoln{0}
% Some commands to allow solutions to be embedded in the assignment file.
\ifcsname issoln\endcsname \else \def\issoln{0} \fi
% Default to an empty solutions environ.
\NewEnviron{soln}{}{}
% Default to an empty qauthor environ.
\NewEnviron{qauthor}{}{}
% Default to visible (but empty) solution box.
\newtcolorbox[]{studentsolution}[1][]{%
    breakable,
    enhanced,
    colback=white,
    title=Solution,
    #1
}

\if\issoln 1
% Otherwise, include solutions as below.
\RenewEnviron{soln}{
    \leavevmode\color{red}\ignorespaces
    \textbf{Solution} \BODY
}{}
\fi

\if\issoln 1
% Otherwise, include solutions as below.
\RenewEnviron{solution}{}
\fi

%%%%%%%%%%%%%%%%%%%%%%%%%%%%%%%%%%%%%%%%%%%
% Commands for customizing the assignment %
%%%%%%%%%%%%%%%%%%%%%%%%%%%%%%%%%%%%%%%%%%%

\newcommand{\courseNum}{\href{https://geometric3d.github.io}{16822}}
\newcommand{\courseName}{\href{https://geometric3d.github.io}{Geometry-based Methods in Vision}}
\newcommand{\courseSem}{\href{https://geometric3d.github.io}{Fall 2022}}
\newcommand{\courseUrl}{\url{https://piazza.com/cmu/fall2022/16822}}
\newcommand{\hwNum}{Problem Set 1}
\newcommand{\hwTopic}{Linear Algebra }
\newcommand{\hwName}{\hwNum: \hwTopic}
\newcommand{\outDate}{Sep. 06, 2022}
\newcommand{\dueDate}{Sep. 13, 2022 11:59 PM}
\newcommand{\instructorName}{Shubham Tulsiani}
\newcommand{\taNames}{Mosam Dabhi, Kangle Deng, Jenny Nan}

%\pagestyle{fancyplain}
\lhead{\hwName}
\rhead{\courseNum}
\cfoot{\thepage{} of \numpages{}}

\title{\textsc{\hwName}} % Title


\author{}

\date{}

%%%%%%%%%%%%%%%%%%%%%%%%%%%%%%%%%%%%%%%%%%%%%%%%%
% Useful commands for typesetting the questions %
%%%%%%%%%%%%%%%%%%%%%%%%%%%%%%%%%%%%%%%%%%%%%%%%%

\newcommand \expect {\mathbb{E}}
\newcommand \mle [1]{{\hat #1}^{\rm MLE}}
\newcommand \map [1]{{\hat #1}^{\rm MAP}}
\newcommand \argmax {\operatorname*{argmax}}
\newcommand \argmin {\operatorname*{argmin}}
\newcommand \code [1]{{\tt #1}}
\newcommand \datacount [1]{\#\{#1\}}
\newcommand \ind [1]{\mathbb{I}\{#1\}}

\newcommand{\blackcircle}{\tikz\draw[black,fill=black] (0,0) circle (1ex);}
\renewcommand{\circle}{\tikz\draw[black] (0,0) circle (1ex);}

\newcommand{\pts}[1]{\textbf{[#1 pts]}}

%%%%%%%%%%%%%%%%%%%%%%%%%%
% Document configuration %
%%%%%%%%%%%%%%%%%%%%%%%%%%

% Don't display a date in the title and remove the white space
\predate{}
\postdate{}
\date{}

% Don't display an author and remove the white space
%\preauthor{}
%\postauthor{}

%%%%%%%%%%%%%%%%%%
% Begin Document %
%%%%%%%%%%%%%%%%%%


\begin{document}

\section*{}
\begin{center}
  \textsc{\LARGE \hwNum} \\
  \textsc{\LARGE \hwTopic\footnote{Compiled on \today{} at \currenttime{}}} \\
  \vspace{1em}
  \textsc{\large \courseNum{} \courseName{} (\courseSem)} \\
  %\vspace{0.25em}
  \courseUrl\\
  \vspace{1em}
  OUT: \outDate \\
  %\vspace{0.5em}
  DUE: \dueDate \\
  Instructor: \instructorName \\
  TAs: \taNames
\end{center}

\section*{START HERE: Instructions}
\begin{itemize}
\item \textbf{Collaboration policy:} All are encouraged to work together BUT you must do your own work (code and write up). If you work with someone, please include their name in your write up and cite any code that has been discussed. If we find highly identical write-ups or code without proper accreditation of collaborators, we will take action according to university policies, i.e. you will likely fail the course. See the \href{https://www.dropbox.com/s/z6o0tinc9eaez46/L01_Overview.pdf?dl=0}{Academic Integrity Section} detailed in the initial lecture for more information.

\item\textbf{Late Submission Policy:} There are \textbf{no} late days for Problem Set submissions.

\item\textbf{Submitting your work:}

\begin{itemize}

% Since we are not using Canvas this semester.
% \item \textbf{Canvas:} We will use an online system called Canvas for short answer and multiple choice questions. You can log in with your Andrew ID and password. (As a reminder, never enter your Andrew password into any website unless you have first checked that the URL starts with "https://" and the domain name ends in ".cmu.edu" -- but in this case it's OK since both conditions are met).  You may only \textbf{submit once} on canvas, so be sure of your answers before you submit.  However, canvas allows you to work on your answers and then close out of the page and it will save your progress.  You will not be granted additional submissions, so please be confident of your solutions when you are submitting your assignment.

\item We will be using Gradescope (\url{https://gradescope.com/}) to submit the Problem Sets. Please use the provided template. Submissions can be written in LaTeX. Regrade requests can be made, however this gives the TA the opportunity to regrade your entire paper, meaning if additional mistakes are found then points will be deducted.
Each derivation/proof should be  completed on a separate page. For short answer questions you \textbf{should} include your work in your solution.  %{\color{red} Since this assignment is for self-review. For this assignment only, you are not required to submit the solution for grading. Solution for this homework will be released on Piazza.}

\end{itemize}

\item \textbf{Materials:} The data that you will need in order to complete this assignment is posted along with the writeup and template on Piazza.

\end{itemize}

For multiple choice or select all that apply questions, replace \lstinline{\choice} with \lstinline{\CorrectChoice} to obtain a shaded box/circle, and don't change anything else.

\clearpage

\section*{Instructions for Specific Problem Types}

For ``Select One" questions, please fill in the appropriate bubble completely:

\begin{quote}
\textbf{Select One:} Who taught this course?
     \begin{checkboxes}
     \CorrectChoice Shubham Tulsiani
     \choice Marie Curie
     \choice Noam Chomsky
     \choice Martial Hebert
    \end{checkboxes}
\end{quote}

For ``Select all that apply" questions, please fill in all appropriate squares completely:

\begin{quote}
\textbf{Select all that apply:} Which are scientists?
{
    \checkboxchar{$\Box$} \checkedchar{$\blacksquare$}
    \begin{checkboxes}
     \CorrectChoice Stephen Hawking
     \CorrectChoice Albert Einstein
     \CorrectChoice Isaac Newton
     \choice None of the above
    \end{checkboxes}
    }
\end{quote}

For questions where you must fill in a blank, please make sure your final answer is fully included in the given space. You may cross out answers or parts of answers, but the final answer must still be within the given space.

\begin{quote}
\textbf{Fill in the blank:} What is the course number?

\begin{tcolorbox}[fit,height=1cm, width=4cm, blank, borderline={1pt}{-2pt},nobeforeafter, halign=center, valign=center]
    \begin{center}\huge16-822\end{center}
    \end{tcolorbox}\hspace{2cm}
\end{quote}

% added course policy section
% currently it is stated that these questions will not be graded, but must be finished
\clearpage

% \section{Written Questions}

\section{Vector Spaces  [8pts]}
\begin{questions}

\question \textbf{[2 pts]} Which of the following subsets of $\mathbb{R}^3$ are vector spaces?
    \textbf{Select all that are true:} \textbf{[2 pts]}
    \checkboxchar{$\Box$} \checkedchar{$\blacksquare$}
    \begin{checkboxes}
        \choice The plane formed by the vector $(v_1, v_2, v_3) $ such that $v_1 = v_2$
        \choice The plane formed by the vector $(v_1, v_2, v_3) $ such that $v_1 = 1$
        \choice The plane formed by the vector $(v_1, v_2, v_3) $  such that $v_1 v_2 v_3 = 0$
        \choice All linear combinations of $\vv = (1, 4, 0)$ and $\wv = (2, 2, 2)$
    \end{checkboxes}

\question \textbf{[2 pts]} For the following questions, consider a matrix $\Av \in \mathbb{R}^{m \times n}$ and a vector $b \in \mathbb{R}^n$. Answer True or false (with a counterexample if false):
    \textbf{Select all that are true:} 
    \begin{checkboxes}
        \choice The vectors b that are not in the column space $\Cv(\Av)$ form a subspace.
        \choice If $\Cv(\Av)$ contains only zero vectors, then $\Av$ is the zero matrix.
        \choice The column space of the matrix $2\Av$ equals the column space of $\Av$
        \choice The column space of the matrix $\Av - \Iv$ equals the column space of $\Av$
    \end{checkboxes}

\question \textbf{[2 pts]} Create a $3 \times 4$ matrix whose solution to $\Av \xv = 0$ is the $\sv_1=\begin{bmatrix}
         -3 \\
         1 \\
         0 \\
         0
        \end{bmatrix}$ and $\sv_2=\begin{bmatrix}
         -2 \\
         0 \\
         -6 \\
         0
        \end{bmatrix}$

    \begin{tcolorbox}[fit,height=5cm, width=\textwidth, blank, borderline={0.5pt}{-2pt},halign=center, valign=center, nobeforeafter]
        %
    \end{tcolorbox}

% \question Answer True or False for the following questions. \textbf{Select all that are True}

%     \textbf{Select all that apply:}
%     \begin{checkboxes}
%     {%
%     \checkboxchar{$\Box$} \checkedchar{$\blacksquare$}
%         \choice A square matrix has no free variables.
%         \choice An invertible matrix has no free variables.
%         \choice A $m \times n$ matrix has more than n pivot variables.
%         \choice A $m \times n$ matrix has no more than m pivot variables.
%         }
%     \end{checkboxes}
% \clearpage
% \question For the given matrix $$\Av = \begin{bmatrix}
%          1 & 0 & 0 \\
%          2 & 1 & 0 \\
%          5 & 0 & 1
%         \end{bmatrix}
% 	\begin{bmatrix}
%          1 & 3 & 0 & 5 \\
%          0 & 0 & 1 & 6 \\
%          0 & 0 & 0 & 0
%         \end{bmatrix}$$
%         Without explicitly computing $\Av$ find the bases and dimensions of all four fundamental subspaces of $\Av$:

%     \begin{tcolorbox}[fit,height=5cm, width=\textwidth, blank, borderline={0.5pt}{-2pt},halign=center, valign=center, nobeforeafter]
%         %
%     \end{tcolorbox}

% \question If a $7{\times}9$ matrix has rank $5$, what are the dimensions of the four subspaces? What is the sum of all four dimensions?

%     \begin{tcolorbox}[fit,height=3cm, width=\textwidth, blank, borderline={0.5pt}{-2pt},halign=center, valign=center, nobeforeafter]
%         %
%     \end{tcolorbox}

\question \textbf{[2 pts]} If a $3{\times}4$ matrix has rank $3$, what are its column space and left nullspace?

    \begin{tcolorbox}[fit,height=3cm, width=\textwidth, blank, borderline={0.5pt}{-2pt},halign=center, valign=center, nobeforeafter]
        %
    \end{tcolorbox}

% \question Determine the value of $x$ for which $\begin{bmatrix}
%     1 \\ x \\ 2
% \end{bmatrix}$ is in the span of $\mathcal{S} = \Biggl\{ \begin{bmatrix}
%     1 \\ 2 \\ -1
% \end{bmatrix} , \begin{bmatrix}
%     -1 \\ -2 \\ 2 
% \end{bmatrix} \Biggl\} $

%     \begin{tcolorbox}[fit,height=3cm, width=\textwidth, blank, borderline={0.5pt}{-2pt},halign=center, valign=center, nobeforeafter]
%         %
%     \end{tcolorbox}

\end{questions}
\clearpage
\section{Eigenvalues, Eigenvector, Singular Value Decomposition [16 pts]}
\begin{questions}
    \question \textbf{[2 pts]}  Deduce the Eigenvalue and Eigenvectors of $\Av$:
    \begin{align*}
        \begin{bmatrix}
            1 & 2 \\ 2 & 4
        \end{bmatrix}
    \end{align*}
    \begin{tcolorbox}[fit,height=3cm, width=\textwidth, blank, borderline={0.5pt}{-2pt},halign=center, valign=center, nobeforeafter]
        %
    \end{tcolorbox}

    \question \textbf{[2 pts]}   For
    \begin{align*}
         \Av = \begin{bmatrix}
             2 & -1 \\ -1 & 2
         \end{bmatrix}
     \end{align*}
     Find the Eigenvalues and Eigenvectors of $\Av, \Av^2$ and $\Av^{-1}$ and $\Av + 4\Iv$

    \begin{tcolorbox}[fit,height=6cm, width=\textwidth, blank, borderline={0.5pt}{-2pt},halign=center, valign=center, nobeforeafter]
    %
    \end{tcolorbox}

    \question \textbf{[2 pts]}  $\Av \in \mathbb{R}^{m \times n}$ is positive definite if for any \textbf{non-zero} vector $\xv \in \mathbb{R}^{n}$  we have $\xv^\top \Av \xv > 0$:
    \begin{align*}
        \xv^\top \Av \xv = \begin{bmatrix}
            x & y
        \end{bmatrix} \begin{bmatrix}
            a & b \\ b & c
        \end{bmatrix}
        \begin{bmatrix}
            x \\ y
        \end{bmatrix}
    \end{align*}

    Test matrices $\Cv$ and $\Dv$ for positive definitiveness
    \begin{align*}
        \Cv=\begin{bmatrix}
                 2 & -1 & 0 \\
                 -1 & 2 & -1 \\
                 0 & -1 & 2
                \end{bmatrix};
        \Dv=\begin{bmatrix}
                 2 & -1 & b \\
                 -1 & 2 & -1 \\
                 b & -1 & 2
                \end{bmatrix}
    \end{align*}
    \begin{tcolorbox}[fit,height=3cm, width=\textwidth, blank, borderline={0.5pt}{-2pt},halign=center, valign=center, nobeforeafter]
    %
    \end{tcolorbox}

    % \clearpage
    % \question When is the symmetric block matrix $\mathbf{M}$ positive definite:
    % \begin{align*}
    %     M=\begin{bmatrix}
    %              A &  B  \\
    %              B^{T} & C
    %             \end{bmatrix}
    % \end{align*}
    % \begin{tcolorbox}[fit,height=3cm, width=\textwidth, blank, borderline={0.5pt}{-2pt},halign=center, valign=center, nobeforeafter]
    % %
    % \end{tcolorbox}

    % \question Without any equation, find the SVD of $\Av$
    % \begin{align*}
    %     \Av = \begin{bmatrix}
    %         2 & 2 \\ 1 & 1
    %     \end{bmatrix}
    % \end{align*}

    % \begin{tcolorbox}[fit,height=3cm, width=\textwidth, blank, borderline={0.5pt}{-2pt},halign=center, valign=center, nobeforeafter]
    % %
    % \end{tcolorbox}

    \question \textbf{[3 pts]}   Estimate the singular values $\sigma_1$ and $\sigma_2$ of the matrix $\Av$
    \begin{align*}
        \Av = \begin{bmatrix}
            1 & 0 \\ C & 1
        \end{bmatrix}
    \end{align*}

    \begin{tcolorbox}[fit,height=2cm, width=\textwidth, blank, borderline={0.5pt}{-2pt},halign=center, valign=center, nobeforeafter]
    %
    \end{tcolorbox}

    \question \textbf{[3 pts]}  Find the pseudoinverse of $\Av \in \mathbb{R}^{m \times n}$
    \begin{align*}
        \Av = \begin{bmatrix}
            2 & 2 \\ 1 & 1
        \end{bmatrix}
    \end{align*}
    \begin{tcolorbox}[fit,height=2cm, width=\textwidth, blank, borderline={0.5pt}{-2pt},halign=center, valign=center, nobeforeafter]
    %
    \end{tcolorbox}

    % \question Find the inverse and the eigenvalues and the determinant of $\Av$:
    % \begin{align*}
    %     A=\begin{bmatrix}
    %          4 & -1 & -1 & 1  \\
    %         -1 & 4 & -1 & -1 \\
    %         -1 & -1 & 4 & -1 \\
    %         -1 & -1 & -1 & 4\\
    %     \end{bmatrix}
    % \end{align*}
    % Describe an Eigenvector matrix $\Sv$ such that $\Sv^{-1} \Av \Sv = \Lambdav$

    % \begin{tcolorbox}[fit,height=5cm, width=\textwidth, blank, borderline={0.5pt}{-2pt},halign=center, valign=center, nobeforeafter]
    % %
    % \end{tcolorbox}
    
    \question \textbf{[4 pts]} Suppose the following information is known about matrix $\mathbf{A}$:
    
    \begin{align*}
    \mathbf{A} \begin{bmatrix}
        1 \\2 \\1
    \end{bmatrix} = 6\begin{bmatrix}
        1 \\ 2 \\1
    \end{bmatrix}, \quad \mathbf{A}\begin{bmatrix}
        1 \\ -1 \\ 1
    \end{bmatrix} = 3 \begin{bmatrix}
        1 \\ -1 \\ 1
    \end{bmatrix}, \quad \mathbf{A} \begin{bmatrix}
        2 \\ -1 \\ 0
    \end{bmatrix} = 3 \begin{bmatrix}
        1 \\ -1 \\ 1
    \end{bmatrix}
    \end{align*}
    \begin{enumerate}[label=\Roman*]
        \item \textbf{[2 pts]} Find the eigenvalues of $\mathbf{A}$
        % \item Find the corresponding eigenspaces
        \item \textbf{[3 pts]} In each of the following subquestions, please justify with a reason (based on the theory of eigenvalues and eigenvectors).
        \begin{enumerate}
            \item Is $\mathbf{A}$ a diagonalizable matrix?
            \item Is $\mathbf{A}$ an invertible matrix?
            % \item Is $\mathbf{A}$ a projection matrix?
        \end{enumerate}

    \end{enumerate}
    \begin{tcolorbox}[fit,height=3cm, width=\textwidth, blank, borderline={0.5pt}{-2pt},halign=center, valign=center, nobeforeafter]
    %
    \end{tcolorbox}     
    
\end{questions}

\clearpage
% \subsection{Planar Homographies: Theory}


% Suppose we have two cameras $\Cv_1$ and $\Cv_2$ looking at a common plane $\Piv$ in $\mathbb{R}^3$.
% Any point $\Pv \in \mathbb{R}^3$ on $\Piv$ projects to a point $\pv \equiv \begin{bmatrix} u_1 & v_1 & 1 \end{bmatrix}^\top \in \mathbb{P}^1$ in camera $\Cv_1$ and
% $\qv \equiv \begin{bmatrix} u_2 & v_2 & 1 \end{bmatrix}^\top \in \mathbb{P}^1$ in camera$\Cv_2$.
% Since $\Pv$ lies on plane $\Piv$, there exists a relationship between $\pv$ and $\qv$.

% In particular, there exists a common $3 \times 3$ matrix $\Hv$, so that for any $\pv$ and $\qv$, the following conditions holds:
% \begin{equation}
% \mathbf{p \equiv Hq} \label{eq:homography}
% \end{equation}

% We call this relationship \emph{planar homography}.

% \textbf{Note:} that both $\pv$ and $\qv$ are in homogeneous coordinates and the equivalence $\equiv$ means $\pv \propto \mathbf{Hq}$ (recall homogeneous
% coordinates). It turns out in fact that this relationship is also true for cameras that are related by pure rotation but not looking points that lie on a plane.

% \begin{questions}
%     \question  Prove that there exists an $\Hv$ that satisfies \autoref{eq:homography} given two camera projection matrices $\Mv_1 \in \Rb^{3 \times 4}$ and $\Mv_2 \in \Rb^{3 \times 4}$ corresponding to cameras
%     $Cv_1$, $\Cv_2$ respectively and a plane $\Piv$.  Produce only the proof of existence of $\Hv$ and not the algebraic expression. Do not produce an actual algebraic expression for $\mathbf{H}$.

%     \textbf{Note:} A degenerate case may happen when the
%     plane $\Pi$ contains both cameras' centers, in which case there are infinite
%     choices of $\mathbf{H}$ satisfying \autoref{eq:homography}.  You can
%     ignore this case in your answer.

%     \begin{tcolorbox}[fit,height=5cm, width=\textwidth, blank, borderline={0.5pt}{-2pt},halign=center, valign=center, nobeforeafter]
%     %
%     \end{tcolorbox}

%     \question Prove that there exists an $\Hv$ that satisfies \autoref{eq:homography} given two cameras separated by a pure rotation.
%     i.e., For $\Cv_1$
%     \begin{align*}
%         \pv = \Kv_1 \begin{bmatrix} \Iv & 0 \end{bmatrix} \Pv
%     \end{align*}
%     For $\Cv_2$
%     \begin{align*}
%         \qv = \Kv_2 \begin{bmatrix} \Rv & 0 \end{bmatrix} \Pv
%     \end{align*}
%     Note that $\Kv$ is not the same for the two cameras.

%     \begin{tcolorbox}[fit,height=5cm, width=\textwidth, blank, borderline={0.5pt}{-2pt},halign=center, valign=center, nobeforeafter]
%     %
%     \end{tcolorbox}

%     \question Why is a planar homography not sufficient to map an arbitrary image of scene to another viewpoint?  State your answer concisely in one or
%      two sentences.

%     \begin{tcolorbox}[fit,height=3cm, width=\textwidth, blank, borderline={0.5pt}{-2pt},halign=center, valign=center, nobeforeafter]
%     %
%     \end{tcolorbox}

%     \textbf{\underline Use the following information to answer questions 4-9}
%     We have a set of points $\pv = \{p_1, p_2, p_3, \dots, p_N\}$ in an image taken by camera $\Cv_1$ and corresponding points $\qv = \{q_1, q_2, q_3, \dots, q_N\}$ in an image taken by $\Cv_2$.
%     Suppose we know that there exists an unknown homography $\Hv$ between corresponding points $\forall i \in \{1, 2, \dots, N\}$. Formally
%     \begin{align*}
%         \Hv : \Pb^1 \rightarrow \Pb^1 ;~~ p_i \equiv \Hv q_i,~~\forall i \in \{1, 2, \dots, N\} \label{eq:2}
%     \end{align*}
%     \question Given $N$ correspondences in $\pv$ and $\qv$ and using \label{eq:2}, derive a set of 2N independent linear equations in the form:
%     \begin{align*}
%         \Av \hv = \mathbf{0}
%     \end{align*}
%     where $\hv$ is a vector of the elements of $\Hv$ and $\Av$ is a matrix composed of elements derived from the point coordinates. Write out an expression for $\Av$.

%     \begin{tcolorbox}[fit,height=4cm, width=\textwidth, blank, borderline={0.5pt}{-2pt},halign=center, valign=center, nobeforeafter]
%     %
%     \end{tcolorbox}

%     \clearpage

%     \question How many elements are there in $\hv$?

%     \begin{tcolorbox}[fit,height=1cm, width=3cm, blank, borderline={0.5pt}{-2pt},halign=center, valign=center, nobeforeafter]
%     %
%     \end{tcolorbox}

%     \question Given $N$ point pairs (correspondences), what are the dimensions of $\Av$?

%     \begin{tcolorbox}[fit,height=1cm, width=3cm, blank, borderline={0.5pt}{-2pt},halign=center, valign=center, nobeforeafter]
%     %
%     \end{tcolorbox}

%     \question How many point pairs (correspondences) are required to solve the system?

%     \begin{tcolorbox}[fit,height=1cm, width=3cm, blank, borderline={0.5pt}{-2pt},halign=center, valign=center, nobeforeafter]
%     %
%     \end{tcolorbox}

%     \question Show how to estimate the elements in $\hv$ to find a solution to minimize this homogeneous linear least squares system.

%     \begin{tcolorbox}[fit,height=4cm, width=\textwidth, blank, borderline={0.5pt}{-2pt},halign=center, valign=center, nobeforeafter]
%     %
%     \end{tcolorbox}

% \end{questions}

% \clearpage
% \textbf{Collaboration Questions} Please answer the following:

% \begin{enumerate}
%     \item Did you receive any help whatsoever from anyone in solving this assignment?
%     \begin{checkboxes}
%      \choice Yes
%      \choice No
%     \end{checkboxes}
%     \begin{itemize}
%         \item If you answered `yes', give full details:
%         \item (e.g. “Jane Doe explained to me what is asked in Question 3.4”)
%     \end{itemize}

%     \begin{tcolorbox}[fit,height=3cm,blank, borderline={1pt}{-2pt},nobeforeafter]
%     %Input your solution here.  Do not change any of the specifications of this solution box.
%     \end{tcolorbox}

%     \item Did you give any help whatsoever to anyone in solving this assignment?
%     \begin{checkboxes}
%      \choice Yes
%      \choice No
%     \end{checkboxes}
%     \begin{itemize}
%         \item If you answered `yes', give full details: \underline{No}
%         \item (e.g. “I pointed Joe Smith to section 2.3 since he didn’t know how to proceed with Question 2”)
%     \end{itemize}

%     \begin{tcolorbox}[fit,height=3cm,blank, borderline={1pt}{-2pt},nobeforeafter]
%     %Input your solution here.  Do not change any of the specifications of this solution box.
%     \end{tcolorbox}

%     \item Did you find or come across code that implements any part of this assignment ? \\Yes/No. (See below policy on “found code”)
%     \begin{checkboxes}
%      \choice Yes
%      \choice No
%     \end{checkboxes}
%     \begin{itemize}
%         \item If you answered `yes', give full details: \underline{No}
%         \item (book \& page, URL \& location within the page, etc.).
%     \end{itemize}
%     \begin{tcolorbox}[fit,height=3cm,blank, borderline={1pt}{-2pt},nobeforeafter]
%     %Input your solution here.  Do not change any of the specifications of this solution box.
%     \end{tcolorbox}
% \end{enumerate}

\end{document}
